\section{Few BS about Lorentz Group}\label{sec:appA}
We will discuss about Lorentz group here. The Lorentz group, denoted by $O(1,3)$, is defined to be the group of all transformations (Lorentz transformation or LT) that preserves the Minkowski metric, that is,
\begin{equation}\label{eqn:lgrp}
    \Lambda^\mathrm{T} \eta \Lambda = \eta\ \iff \ \eta_{\mu\nu}\trmat{\mu}{\alpha}\trmat{\nu}{\beta}=\eta_{\alpha\beta}
\end{equation}
Note that, as long as we are not using the indices, we are talking about the \textit{abstract} group element. However, when we consider the expression using index notation, we are essentially considering the matrix representation of the group element using the \textit{standard basis} of the Minkowski space\footnote{The Mindkowski space is denoted by $\mathbb{R}^{1,3}$ which is equal to $\mathbb{R}^4$ along with the associated metric $\eta_{\mu\nu}$. $(1,3)$ refers to the metric signature, that is, 1 positive entry and 3 negative entries viz. $\mathrm{diag}(+1,-1,-1,-1)$}.\\[0.2cm] From the above definition, we can instantly see that for any $\Lambda \in O(1,3)$, $\det\Lambda = \pm 1$. For $\det \Lambda =+1$, we refer to them as \textit{proper LT} (denoted by $L_{+}$) and for $\det \Lambda =-1$ (denote by $L_-$), we refer to them as \textit{improper LT}.\\
In Eqn. \ref{eqn:lgrp}, if we take $\alpha=\beta=0$, we get:
\begin{align*}
   1 =  \eta_{\mu\nu}\trmat{\mu}{0}\trmat{\nu}{0} =\qty( \trmat{0}{0})^2 - \sum\limits_{i=1}^3 \qty( \trmat{i}{0})^2 \implies \qty( \trmat{0}{0})^2 \geq 1
\end{align*}
Thus, we can have $\trmat{0}{0}\geq 1$ which we refer to as \textit{orthochronous} (denoted by $L^{\uparrow}$) or $\trmat{0}{0}\leq -1$ which we call \textit{antichronous} (denoted by $L^{\downarrow}$). \\[0.2cm]
From the above two classifications using the signs of the determinant and $\trmat{0}{0}$, we can thus have four disconnected components of the Lorentz group:
\begin{itemize}
    \item $L^{\uparrow}_{+}$: These are the proper orthochronous LT. The identity $\mathds{1}$ is contained in this subgroup. It is denoted by by $\mathrm{SO}^+(1,3)$ and is continuously connected to the identity, that is, we can build up any
finite Lorentz transformation by making many consecutive small Lorentz transformations. Sometimes, as an abuse of terminology, this subgroup is itself called the \textit{Lorentz group}, since this contains the only physical transformations. 
      \item $L^{\uparrow}_{-}$: These are written as a product of the \textit{parity} and $L^{\uparrow}_{+}$ transformations.
    \item $L^{\downarrow}_{+}$: These contain both time and space inversion and hence are written as a product of $PT$ and $L^{\uparrow}_{+}$ transformation.
    \item $L^{\downarrow}_{-}$: These contain only time inversion and hence are written as a product of \textit{time reversal } and $L^{\uparrow}_{+}$ transformation. 
\end{itemize} 
Thus, the entire Lorentz group can be written symbolically:
$$L = L^{\uparrow}_{+} \bigcup PL^{\uparrow}_{+} \bigcup PT L^{\uparrow}_{+} \bigcup TL^{\uparrow}_{+}$$
Note that except the proper orthochronous LT, none of the other component can form subgroups since these do not contain the identity element.
\subsection{Spacetime Representation}
For the time being, we will consider only the proper, orthochronous Lorentz group, that is, $\mathrm{SO}^+(1,3)$ which is a Lie group. Since this is continuously connected to the identity, let us consider an infinitesimal transformation from the identity given by,
$$\trmat{\mu}{\nu} = \delt{\mu}{\nu}+\tensor{\omega}{^\mu _\nu}$$
Substituting this in Eq. \ref{eqn:lgrp} and retaining terms upto the first oder in $\omega$, we have the following:
\begin{align*}
    \eta_{\alpha \beta} &= \eta_{\mu\nu} \qty(\delt{\mu}{\alpha} + \tensor{\omega}{^\mu _\alpha})\qty(\delt{\nu}{\beta} + \tensor{\omega}{^\nu _\beta})\\
    &=\eta_{\alpha\beta} + \delt{\mu}{\alpha}\eta_{\mu\nu}\tensor{\omega}{^\nu _\beta} + \eta_{\mu\nu}\delt{\nu}{\beta} \tensor{\omega}{^\mu _\alpha}\\
    &=\eta_{\alpha\beta} + \delt{\mu}{\alpha}\tensor{\omega}{_\mu _\beta} + \delt{\nu}{\beta} \tensor{\omega}{_\nu _\alpha}\\
    &=\eta_{\alpha\beta} + \tensor{\omega}{_\alpha _\beta} +  \tensor{\omega}{_\beta _\alpha}
\end{align*}
From this, we get that infinitesimal displacement $\omega$ must be anti-symmetric, that is, 
$$\boxed{\omega_{\alpha\beta} = -\omega_{\beta\alpha}}$$
We know that antisymmetric matrices have six independent elements (since diagonals are zero and upper-diagonal has $(16-4)/2 = 6$ elements and lower diagonal are negative of the upper diagonal) and hence the basis of antisymmetric matrices consists six elements. Of these six, three corresponds to \textit{rotations} while the other three corresponds to \textit{boosts}\\[0.2cm] Let us consider this basis of the second rank anti-symmetric tensor $\{\SJ^{\rho\sigma}\}$ which forms the \textit{generators} of the transformation. These are antisymmetric in the indices $\rho, \sigma$ and hence there are six of them. These can be written as:
$$\tensor{{\qty(\SJ^{\rho\sigma})}}{^\mu _\nu} :=i \qty(\eta^{\mu\rho} \tensor{\delta}{^\sigma _\nu} - \eta^{\mu\sigma} \tensor{\delta}{^\rho _\nu})$$
As $\omega$ is antisymmetric, we can write it as a linear combination of the basis elements, that is, 
$$\tensor{\omega}{^\mu _\nu} =-i \Omega_{\rho\sigma}\tensor{{\qty(\SJ^{\rho\sigma})}}{^\mu _\nu}$$
Then we can write the infinitesimal transformation in the following way:
$$\Lambda = \mathds{1} -i \Omega_{\rho\sigma}\SJ^{\rho\sigma}$$
Any finite Lorentz transformation can be generated by exponentiating the generators, that is, 
$$\Lambda_{\mathrm{finite}} = \exp\qty(-\frac{i}{2}\Omega_{\rho\sigma}\SJ^{\rho\sigma})$$
Once we found the generators, any element of the Lie algebra of the Lorentz group will be a linear combination of them. It is a good time for us to calculate the Lie bracket since it is an essential characteristic of the Lie algebra. Since any element is a linear combination of the generators, it is sufficient to calculate the Lie bracket for the generators only, which we can do by the explicit form that we wrote earlier. 
\begin{align*}
    \tensor{{\comtr{\SJ^{\mu\nu}}{\SJ^{\rho\sigma}}}}{^\alpha _\beta} &=  \tensor{\qty({{\SJ^{\mu\nu}}{\SJ^{\rho\sigma}}})}{^\alpha _\beta} - \tensor{\qty({{\SJ^{\rho\sigma}}{\SJ^{\mu\nu}}})}{^\alpha _\beta}\\
    &=\tensor{(\SJ^{\mu\nu})}{^\alpha _\theta}\tensor{(\SJ^{\rho\sigma})}{^\theta _\beta} -\tensor{(\SJ^{\rho\sigma})}{^\alpha _\theta}\tensor{(\SJ^{\mu\nu})}{^\theta _\beta}\\
    &=( \eta^{\mu\alpha} \tensor{\delta}{^\nu _\theta} - \eta^{\nu\alpha} \tensor{\delta}{^\mu _\theta})( \eta^{\rho\theta} \tensor{\delta}{^\sigma _\beta} - \eta^{\sigma\theta} \tensor{\delta}{^\rho _\beta}) -  ( \eta^{\rho\alpha} \tensor{\delta}{^\sigma _\theta} - \eta^{\sigma\alpha} \tensor{\delta}{^\rho _\theta}) ( \eta^{\mu\theta} \tensor{\delta}{^\nu _\beta} - \eta^{\nu\theta} \tensor{\delta}{^\mu _\beta})\\
    &=\qty{\eta^{\mu\alpha} \tensor{\delta}{^\nu _\theta} \eta^{\rho\theta} \tensor{\delta}{^\sigma _\beta}  - \eta^{\mu\alpha} \tensor{\delta}{^\nu _\theta}\eta^{\sigma\theta} \tensor{\delta}{^\rho _\beta} -\eta^{\nu\alpha} \tensor{\delta}{^\mu _\theta}\eta^{\rho\theta} \tensor{\delta}{^\sigma _\beta} +\eta^{\nu\alpha} \tensor{\delta}{^\mu _\theta}\eta^{\sigma\theta} \tensor{\delta}{^\rho _\beta}} \\
    & \qquad - \qty{\eta^{\rho\alpha} \tensor{\delta}{^\sigma _\theta}\eta^{\mu\theta} \tensor{\delta}{^\nu _\beta}    -  \eta^{\rho\alpha} \tensor{\delta}{^\sigma _\theta}\eta^{\nu\theta} \tensor{\delta}{^\mu _\beta} -\eta^{\sigma\alpha} \tensor{\delta}{^\rho _\theta}\eta^{\mu\theta} \tensor{\delta}{^\nu _\beta}    + \eta^{\sigma\alpha} \tensor{\delta}{^\rho _\theta}\eta^{\nu\theta} \tensor{\delta}{^\mu _\beta}} \\
    &=\qty{\eta^{\mu\alpha}  \eta^{\rho\nu} \tensor{\delta}{^\sigma _\beta}  - \eta^{\mu\alpha} \eta^{\sigma\nu} \tensor{\delta}{^\rho _\beta} -\eta^{\nu\alpha} \eta^{\rho\mu} \tensor{\delta}{^\sigma _\beta} +\eta^{\nu\alpha}\eta^{\sigma\mu} \tensor{\delta}{^\rho _\beta}} \\ &\qquad - \qty{\eta^{\rho\alpha} \eta^{\mu\sigma} \tensor{\delta}{^\nu _\beta}    -  \eta^{\rho\alpha} \eta^{\nu\sigma} \tensor{\delta}{^\mu _\beta} -\eta^{\sigma\alpha}\eta^{\mu\rho} \tensor{\delta}{^\nu _\beta}    + \eta^{\sigma\alpha} \eta^{\nu\rho} \tensor{\delta}{^\mu _\beta}} \textcolor{red}{\text{complete proof!}}
\end{align*}
From this, we find out the Lie bracket between the generators of the Lie Algebra of $\mathrm{SO}(1,3)$:
$$\comtr{\SJ^{\mu\nu}}{\SJ^{\rho\sigma}} =i \qty(\eta^{\nu\rho}\SJ^{\mu\sigma}-\eta^{\mu\rho}\SJ^{\nu\sigma}-\eta^{\nu\sigma}\SJ^{\mu\rho}+\eta^{\mu\sigma}\SJ^{\nu\rho})$$
We had obtained the six generators for the Lie algebra and how let us rearrange these into two separate parts:
$$J^i = \frac{i}{2}\varepsilon^{ijk}\SJ^{jk}\qquad K^i = i\SJ^{i0}$$
In terms of these newly defined quantitites, the Lie bracket becomes \textcolor{red}{check!}:
\begin{align*}
    \comtr{J^i}{J^j}&=i \vep^{ijk}J^k\\
    \comtr{K^i}{K^j}&=-i \vep^{ijk}J^k\\
    \comtr{J^i}{K^j}&=i \vep^{ijk}K^k
\end{align*}
This shows that $J^i$ as defined above kinda denotes angular momentum since it satisfies the Lie Algebra of $\mathrm{SU}(2)$ while the other relation shows that $K^i$ is a spatial vector. Now let us define two quantitites:
$$\theta^i = \frac{1}{2}\vep^{ijk}\Omega_{jk}\qquad \eta^i = \Omega^{i0}$$
Using this we can write the generators as:
$$\frac{1}{2}\Omega_{\mu\nu}\SJ^{\mu\nu} = \Omega_{12}\SJ^{12} + \Omega_{13}\SJ^{13}+ \Omega_{23}\SJ^{23} +  \Omega_{i0}\SJ^{i0} = \bs{\theta}\cdot \vb{J} - \bs{\eta}\cdot \vb{K}$$
Thus any Lorentz transformation can then be written as:
$$\boxed{\Lambda = \exp\qty(-i\ \bs{\theta}\cdot \vb{J} +i\ \bs{\eta}\cdot \vb{K})}$$